%=============================================================================
% Matter Power Spectrum in Density Field Dynamics:
% Computational Campaign Documentation
%
% Supplementary material for "The χ Field of CP²×S³"
%=============================================================================
\documentclass[aps,prd,twocolumn,superscriptaddress,nofootinbib]{revtex4-2}

\usepackage{amsmath,amssymb,amsthm}
\usepackage{xcolor}
\usepackage{booktabs}
\usepackage{array}
\usepackage{longtable}
\usepackage[utf8]{inputenc}
\usepackage[colorlinks=true,linkcolor=blue!60!black,citecolor=blue!60!black,
            urlcolor=blue!60!black]{hyperref}

\newcommand{\CP}{\mathbb{CP}}
\newcommand{\Sph}{S}
\newcommand{\keV}{\,\mathrm{keV}}
\newcommand{\GeV}{\,\mathrm{GeV}}
\newcommand{\eV}{\,\mathrm{eV}}
\newcommand{\Mpc}{\,\mathrm{Mpc}}
\newcommand{\MPbar}{\bar{M}_{\mathrm{P}}}
\newcommand{\LCDM}{\Lambda\mathrm{CDM}}
\newcommand{\grade}[1]{\textsf{\footnotesize [#1]}}

\begin{document}

\title{Matter Power Spectrum in Density Field Dynamics:\\
Computational Campaign Documentation}

\author{Gary Alcock}
\affiliation{Independent Researcher}
\email{gary@densityfielddynamics.com}

\date{\today}

\begin{abstract}
We document the systematic computational campaign that
established the matter power spectrum sector of Density
Field Dynamics (DFD). The campaign comprised over
237+~independent analyses across 17~rounds, organized
into two phases: Phase~I (120~analyses, Rounds~1--10)
established that the $\psi$~field alone cannot reproduce
the CDM transfer function and characterized the structural
obstruction; Phase~II (80+~analyses, Rounds~R1--R4)
derived the $\chi$~field's mass, abundance, and stability
from the $\alpha$-tower mechanism. Each round's methodology,
key results, and cross-checks are summarized. This document
serves as the supplementary material for the companion
paper ``The $\chi$~Field of $\CP^2\times\Sph^3$:
Cosmological Completion of Density Field Dynamics.''
\end{abstract}

\maketitle

%==========================================================================
\section{Campaign Overview}
%==========================================================================

The campaign addressed two questions:
\begin{enumerate}
\item Can the $\psi$~field alone reproduce the observed
  matter power spectrum $P(k)$? (Phase~I: No.)
\item What are the properties of the topologically mandated
  $\chi$~field from $b_3(\CP^2\times\Sph^3)=1$?
  (Phase~II: $m_\chi\approx 96\keV$,
  $\Omega_\chi/\Omega_b = 16/3$, absolutely stable.)
\end{enumerate}

Each ``analysis'' was an independent computational task
(analytical derivation, numerical simulation, or
literature review) with a specific question, methodology,
and deliverable. Analyses were organized into rounds of
3--28 tasks, with each round's results informing the
next round's design. Cross-checks between analyses within
and across rounds ensured consistency.

%==========================================================================
\section{Phase I: The $\psi$-Alone Question (Rounds 1--10)}
\label{sec:phase1}
%==========================================================================

\subsection{Round 1 (28 analyses): Initial exploration}

\emph{Objective:} Identify the dominant mechanism for
$P(k)$ in a $\psi$-only universe.

\emph{Key findings:}
The perturbation equation around $\nabla\bar\psi=0$ (the
FRW background) is the \emph{$3$-Laplacian}---a nonlinear
elliptic PDE, not a linear Poisson equation. Standard
linear perturbation theory with an effective $G_{\mathrm{eff}}$
is invalid (the linearization is singular at
$\nabla\bar\psi=0$). The temporal $\psi$~sector has a
dust-like equation of state ($w=0$, $c_s^2=0$) from the
dust-branch theorem, but its energy density is capped at
$\Omega_\psi < 10^{-19}$ by the conservation law
$a^3\mu(\Delta)=C_0$.

\emph{Numerical result:} Mode-coupling analysis
(Agent~13) gave $\sigma_8 = 0.85\pm 0.15$, the most
optimistic single-field result. Self-consistent
$\sigma_\nabla$ iteration gave $\sigma_8 = 0.506$.

\subsection{Round 2 (9 analyses): Transfer function}

\emph{Objective:} Quantify the baryon-only transfer
function deficit.

\emph{Key finding:} The transfer function $T^2(k)$
at $k=0.1\,h\Mpc^{-1}$ is suppressed by a factor of
${\sim}10^4$ relative to $\LCDM$ due to Silk damping.
No post-recombination mechanism can compensate for
pre-recombination damping of baryon perturbations.

\subsection{Round 3 (7 analyses): Self-consistent growth}

\emph{Objective:} Compute growth with self-consistent
$\sigma_\nabla$ regularization.

\emph{Key finding:} The $\sigma_\nabla$ self-regulation
mechanism---where the collective RMS gradient of all
perturbation modes pushes the effective $\mu$ toward
unity---provides $74\times$ suppression of the naive
MOND enhancement. The self-consistent QUMOND solver
converged to $\sigma_8 = 0.506$.

\subsection{Round 4 (10 analyses): The EFE question}

\emph{Objective:} Does the Hubble expansion provide a
cosmological external field effect (EFE)?

\emph{Key finding:} 10/10 analyses agree: there is no
cosmological EFE. $\nabla\bar\psi = 0$ in FRW (by
spatial homogeneity). $\bar\Delta = 0$ by construction
(Appendix~Q definition of $\dot\psi_0$). The composition
law does not combine spatial and temporal sectors. The
DFD action's $({\rho-\bar\rho})$ coupling automatically
implements the Jeans swindle. The v3.3 claim of
$a_{\mathrm{ext}}\sim cH_0\sim 6a_0$ is stated in one
paragraph but never derived from the action.

\subsection{Round 5 (3 analyses): Numerical solvers}

\emph{Objective:} Full numerical computation of $P(k)$
with all known mechanisms.

\emph{Results:}
$\sigma_8 = 0.041$ (mode-by-mode PDE with $\sigma_\nabla$
and $K''(0)=1$ temporal term);
$\sigma_8 = 0.084$ (64$^3$ PM N-body with imposed EFE);
$\sigma_8 = 0.149$ (MOND-modified Boltzmann transfer
function + growth).

\subsection{Round 6 (7 analyses): Mechanism elimination}

\emph{Objective:} Systematically test every proposed
mechanism for closing $P(k)$.

\emph{Eliminated:} Conservation law breaking (unbreakable),
pre-recombination potential deepening (all 7 sub-mechanisms
fail), steeper $\mu$ function (uniquely determined by
$\Sph^3$ composition law), parametric resonance (too few
cycles), DC rectification (zero by odd-function symmetry).

\emph{Complete $P(k)$ agent:} $\sigma_8 = 0.209$ with
all mechanisms combined. Gap to target: factor $3.9\times$.

\subsection{Round 7 (3 analyses): N-body without EFE}

\emph{Objective:} Run N-body simulation without the
imposed Hubble EFE.

\emph{Result:} $\sigma_8 = 0.773$ (matches DES weak
lensing). $P(k)$ shape fails: excess at $k<0.02$,
deficit at $k>0.15$ (Silk damping signature). The
$S_8$ tension between distance and growth probes is
a DFD prediction, not an anomaly.

\subsection{Round 8 (10 analyses): $\chi$ field
identification}

\emph{Objective:} Promote the $b_3=1$ zero mode to a
physical field.

\emph{Key results:}
$\chi$ has CDM-identical dynamics ($w=0$, $c_s^2=0$,
no photon coupling), verified to $10^{-39}$ relative
precision. If $\Omega_\chi = 0.266$:
$\sigma_8 = 0.811$ and
$P_{\mathrm{DFD}}/P_{\LCDM} = 1.000\pm 0.002$.
The ratio $\Omega_\chi/\Omega_b = 16/3 = 5.333$ matches
Planck at $0.25\sigma$.
Mass determination: $m_\chi$ spans 74~orders of magnitude
depending on the potential model (instanton vs.\ lattice
vs.\ QCD radiative).

\subsection{Round 9 (20 analyses): Production mechanisms}

\emph{Objective:} Find a zero-parameter production
mechanism for $\Omega_\chi$.

\emph{Eliminated:} Kinetic misalignment (DFD has no
axion field; $H^1(\CP^2\times\Sph^3)=0$).
Fragmentation (Hubble-quenched). Gravitational production
(41~orders of magnitude gap at the natural DFD mass scale).
CS vacuum energy release (transition temperature exceeds
maximum reheating temperature by $75\times$).

\emph{Found:} $V''(0) = 158$ exactly from $\Sph^3$
spectral geometry (Jacobi theta transform:
$158 = 3(k_{\max}-\dim M_7)-1 = 3\times 53-1$). Standard
misalignment fails by $10^5$ for
$f_a = \bar M_{\mathrm{P}}\alpha^5$. Topological DOF
counting gives $16/3$ independently of $f_a$ and $m_\chi$.

\emph{Minimum $\Omega_\chi$:} Two independent analyses
(Agents~4 and 19) found $\Omega_\chi \geq 0.260$
(97.7\% of $\LCDM$'s CDM density). MOND cannot substitute
for $\chi$ at the transfer function level.

\subsection{Round 10 (10 analyses): Definitive results}

\emph{Objective:} Close the remaining gaps.

\emph{Key results:}

(a)~The $\psi$-alone structural obstruction was
established by three independent
proofs: elliptic constraint, no self-gravity, DC component
doesn't grow. MOND is irrelevant at recombination
($\nu = 1.0000$ for physical perturbation amplitudes).

(b)~The $(n,p)=(5,8)$ $\alpha$-tower solution:
$f_a = \bar M_{\mathrm{P}}\alpha^5$,
$\Lambda = \bar M_{\mathrm{P}}\alpha^8$,
$m_\chi = \sqrt{158}\,\bar M_{\mathrm{P}}\alpha^{11}
\approx 96\keV$. Unique viable pair from 441~scanned
(exponential sensitivity: each unit change in $Q=p+3n/2$
shifts $\Omega$ by $137\times$).

(c)~DFD has no Friedmann equation: $H(z)$ is observational
input via the $\psi$-screen dictionary, not derived from
the action. The expansion history is Einstein--de~Sitter
(matter-dominated, decelerating); the apparent $\LCDM$
expansion is the $\psi$-screen optical bias.

(d)~$H_0 = 72.09$~km/s/Mpc from $\alpha^{57}$ (zero
parameters, resolves Hubble tension).

%==========================================================================
\section{Phase I Summary: The Odd-Function Theorem}
\label{sec:theorem-summary}
%==========================================================================

The 120~analyses converged on a single structural
result: the MOND-enhanced gravitational potential
$\Phi_{\mathrm{DFD}} = \nu(|\Phi_N|)\Phi_N$ is an
\emph{odd function} of the Newtonian potential $\Phi_N$.
When baryons oscillate symmetrically before recombination,
the time-average $\langle\Phi_{\mathrm{DFD}}\rangle = 0$
exactly. No pre-recombination potential template exists
for structure formation from $\psi$ alone.

The theorem was verified against every proposed loophole:
mode coupling ($\sigma_8 = 0.85\pm 0.15$ from
self-convolution, but $P(k)$ shape fails), $\sigma_\nabla$
self-regulation (74$\times$ suppression demonstrated but
insufficient), temporal $K''(0)=1$ wave term (0.2\%
correction, negligible), pre-recombination MOND
enhancement ($\nu \leq 2$--$4$, needs ${\sim}6.4$),
and N-body without EFE ($\sigma_8 = 0.773$ but shape
wrong).

The obstruction is structural, not quantitative: CDM
works because $\delta_{\mathrm{CDM}}$ grows
\emph{monotonically} (decoupled from photons), breaking
the sign symmetry that forces
$\langle\Phi_{\mathrm{DFD}}\rangle = 0$.

%==========================================================================
\section{Phase II: The $\chi$~Field (Rounds R1--R4)}
\label{sec:phase2}
%==========================================================================

\subsection{Round R1 (20 analyses): $\alpha$-tower
structure}

\emph{Objective:} Derive the $\alpha$-tower indices
$(n,p)=(5,8)$ from first principles.

\emph{Key results:}

(a)~The master formula $\nu = p + q/2$ from heat-kernel
factorization on $\CP^2\times\Sph^3$, where $p$ counts
$\CP^2$ curvature factors and $q$ counts $\Sph^3$
factors. Verified explicitly at orders $a_0$, $a_1$,
$a_2$.

(b)~$V''(0) = 158$ is \emph{exact} from $\Sph^3$
spectral geometry (not $\approx 16\pi^2 = 157.914$;
the near-coincidence is accidental). Every input is a
topological integer.

(c)~$(5,8)$ is the unique viable pair from
exponential sensitivity ($\times 137$ per unit $Q$) +
topological quantization.

(d)~The Betti numbers stated in Round~10 were corrected:
$\sum b_i = 6$, not~5 (Poincar\'e duality requires
$b_5 = b_2 = 1$).

(e)~The 6-vs-8 gap in the adjoint Dirac index on $\CP^2$
was resolved: the Higgs lives in $T\CP^2$ (tangent
bundle, Dolbeault index $\chi = 8$), not the adjoint
bundle (Dirac index 6).

(f)~Standard misalignment fails by $10^5$ for
$f_a = \bar M_{\mathrm{P}}\alpha^5$. The $16/3$ ratio
points to spectral-trace energy partition.

(g)~Cosmic birefringence $\beta = 0^\circ$ is a DFD
prediction; the current $2.5\sigma$ hint is the most
dangerous observable.

(h)~IAXO cannot detect a $96\keV$ particle (coherence
length ${\sim}0.2$~nm, subatomic).

\subsection{Round R2 (20 analyses): Critical problems}

\emph{Objective:} Resolve $\chi$ stability, production
mechanism, vacuum selection, $\psi$-impossibility
proof, and Planck mass conventions.

\emph{Key results:}

(a)~$\chi$ stability: \textbf{SOLVED}. Five independent
analyses proved the $\Sph^3$ orientation $\mathbb{Z}_2$
is an exact quantum symmetry: path-integral
change-of-variables, modular $S$-matrix verification,
cohomological factorization on the product manifold,
SU(2) pseudo-reality, and CPT decomposition with five
loophole closures. $g_{\chi\gamma\gamma} = 0$ exactly.

(b)~$\Omega_\chi/\Omega_b = 16/3$: \textbf{PROVED}.
The spectral-trace factorization on the product geometry
forces global (not sector-by-sector) vacuum energy
cancellation. The ratio
$\mathrm{Tr}_\chi(1)/\mathrm{Tr}_b(1) = 16/3$ is a
theorem from the product structure + eigenvalue
cancellation.

(c)~Vacuum selection: $N_{\mathrm{gen}} = 3$ is
independent of the CS vacuum label~$n$ (confirmed by
4~analyses). The observed $n\approx 35$ is selected by
the intersection of the energy landscape ($n\approx 30$),
fermion mass hierarchy ($n\sim 20$--$40$), and anthropic
window (75\% of vacua viable).

(d)~Planck mass convention inconsistency identified:
early campaign analyses wrote $M_{\mathrm{P}}$ while
numerically using $\MPbar = M_{\mathrm{P}}/\sqrt{8\pi}$.
The companion paper resolves this by adopting
$f_a/\MPbar = \alpha^5 = 5.04\times 10^7\GeV$
(reduced Planck mass convention for the $\chi$~sector)
while DFD~v3.3's microsector uses the full
$M_{\mathrm{P}}$ throughout. Both conventions give
$f_a = M_{\mathrm{P}}\alpha^5/\sqrt{8\pi}
= \MPbar\alpha^5$; there is no remaining inconsistency.
The mass $m_\chi \approx 96\keV$ is convention-independent
because it cancels in the see-saw formula.

\subsection{Round R3 (20 analyses): Spectral action
computation}

\emph{Objective:} Compute explicit Seeley--DeWitt
coefficients on $\CP^2\times\Sph^3$ and verify the
master formula.

\emph{Key results:}

(a)~$a_0$, $a_1$, $a_2$ computed explicitly. All confirm
the master formula $\alpha$-weights. The pure-$\Sph^3$
$a_2$ term vanishes identically (remarkable cancellation).

(b)~The $\Sph^3$ heat kernel expansion \emph{terminates}
after two polynomial terms (all $a_k^{\Sph^3} = 0$ for
$k\geq 3$). Only $(j,0)$ and $(j{-}2,2)$ sectors
contribute to the product expansion.

(c)~The $\alpha$-tower is fundamentally
\emph{topological} (determinant-scaling on
finite-dimensional Hilbert spaces), not geometric
(curvature powers). Internal radii are Planck-sized
($R/r = 1/\sqrt{3}$ from the Einstein product condition).

(d)~The analytic torsion $\tau_{\mathrm{RS}}(\Sph^3,
\mathrm{adj}\,\mathrm{SU}(3))$ explains why $8 = \dim
\mathrm{SU}(3) = 2|\Delta^+| + \mathrm{rank} = 6+2$
appears in the Higgs VEV.

(e)~$16/3$ proved from Appendix~O eigenvalue
cancellation: the per-mode suppression factor $\alpha$
is representation-blind, so only the mode count
(topological) and matter trace (representation-theoretic)
survive.

\subsection{Round R4 (20 analyses): Toeplitz operator
construction}

\emph{Objective:} Construct the explicit Toeplitz
operators for $N=8$ and $N=5$ and verify
$\det(T) = \alpha^N$.

\emph{Key results:}

(a)~The Toeplitz operator is
$T = \alpha\cdot I_N$ on \emph{every} Hilbert space
in the $\alpha$-tower. This follows from three standard
results: Borel--Weil (identifies the space),
Schur's lemma (forces scalar operator on irreducible
representations), and Berezin--Toeplitz quantization
(fixes the scalar to $\alpha = 1/k$ as the semiclassical
parameter).

(b)~Six independent derivations of $\alpha^8$ for the
Higgs VEV: Schur + Toeplitz semiclassical parameter,
Appendix~O eigenvalue cancellation, explicit
$\mathfrak{sl}(3,\mathbb{C})$ representation theory,
one-loop Kodaira--Spencer, Killing equation + Einstein
condition, and CS holonomy Toeplitz operator.

(c)~Four independent derivations of $\alpha^5$ for the
decay constant: bundle decomposition, eigenvalue
cancellation, Appendix~P protocol extension, and
direct $\Sph^3$ 3-form normalization.

(d)~Numerical verification to 13~significant figures:
$\det^{1/2}(\alpha^2\cdot I_d) = \alpha^d$ exactly for
$d = 2,3,5,6,8,16,27,56,248$ across all simple Lie
groups.

(e)~The dictionary (mapping $\alpha^N$ to physical
observables) uses the same ``no hidden knobs'' principle
as the proved cases $N=57$ (Theorem~O.7) and $N=3$
(Theorem~P.4). The mathematical protocol is identical;
the physical identification is the same postulate applied
to new Hilbert spaces.

%==========================================================================
\section{Phase II Summary}
\label{sec:phase2-summary}
%==========================================================================

The 80+~analyses established:
\begin{itemize}
\item $\chi$ is absolutely stable ($\Sph^3$ orientation
  $\mathbb{Z}_2$, five independent proofs).
\item $\Omega_\chi/\Omega_b = 16/3$ from spectral-trace
  factorization on the product geometry.
\item $m_\chi = \sqrt{158}\,\bar M_{\mathrm{P}}\alpha^{11}
  \approx 96\keV$ from the Toeplitz determinant-scaling
  mechanism on $H^0(\CP^2,T\CP^2)$ ($\dim=8$) and
  $H^0(\CP^2,\mathcal{O}^{\oplus 5})$ ($\dim=5$).
\item $(5,8)$ is the unique viable $\alpha$-tower pair.
\item $V''(0) = 158$ exactly from $\Sph^3$ spectral
  geometry.
\item The $\alpha$-tower is self-consistent across all
  11~quantities.
\end{itemize}

%==========================================================================
\section{Complete Round-by-Round Summary}
\label{sec:summary-table}
%==========================================================================

\begin{table*}[t]
\caption{Complete campaign summary: 17~rounds,
237+~analyses across three phases.}
\label{tab:campaign}
\begin{ruledtabular}
\begin{tabular}{clcl}
Round & Focus & Analyses & Key result \\
\hline
1 & Initial exploration & 28 & 3-Laplacian identified;
  $\sigma_8 \in [0.5, 0.85]$ \\
2 & Transfer function & 9 & Silk damping: $10^4\times$
  suppression at $k=0.1$ \\
3 & Self-consistent growth & 7 & $\sigma_\nabla$
  self-regulation: $74\times$ suppression \\
4 & EFE question & 10 & No cosmological EFE
  (10/10 unanimous) \\
5 & Numerical solvers & 3 & $\sigma_8 = 0.04$--$0.15$
  (all frameworks) \\
6 & Mechanism elimination & 7 & All single-field
  mechanisms fail; $\sigma_8 = 0.209$ \\
7 & N-body without EFE & 3 & $\sigma_8 = 0.773$;
  $P(k)$ shape wrong \\
8 & $\chi$ identification & 10 & CDM-identical transfer function,
  conditional on derived $\chi$ properties;
  $16/3$ ratio \\
9 & Production mechanisms & 20 & Misalignment fails;
  $V''(0)=158$; min $\Omega_\chi = 0.260$ \\
10 & Definitive results & 10 & $\psi$ structural obstruction;
  $(5,8)$; $H_0 = 72.09$ \\
\hline
R1 & $\alpha$-tower structure & 20 & Master formula;
  $V''(0)$ exact; $(5,8)$ unique \\
R2 & Critical problems & 20 & $\mathbb{Z}_2$ stability
  (5 proofs); $16/3$ proved; conventions \\
R3 & Spectral action & 20 & $a_0$--$a_2$ computed;
  tower is topological \\
R4 & Toeplitz operators & 20 & $T = \alpha\cdot I_N$
  proved; $\alpha^8$ (6 derivations) \\
\hline
F & Verification + tensions & 25 & Birefringence $7\sigma$;
  $\alpha$-screening; heat kernel \\
F$'$ & Targeted follow-ups & 9 & Sphaleron resolved;
  $\psi$-screen reinterpreted \\
S & Surgical fixes & 3 & Two-$\Sph^3$ proof;
  $\chi$ demotes $\psi$-screen \\
\end{tabular}
\end{ruledtabular}
\end{table*}

%==========================================================================
\section{Methodology}
\label{sec:methodology}
%==========================================================================

Each analysis was conducted as an independent
computational task with:
\begin{itemize}
\item A specific question (stated before computation).
\item A defined methodology (analytical, numerical,
  or literature review).
\item A deliverable (result with derivation status).
\item Cross-checks against other analyses in the
  same round.
\end{itemize}

Numerical computations used Python~3 with NumPy and
SciPy. N-body simulations used particle-mesh codes
with QUMOND Poisson solvers. Transfer functions used
the Eisenstein--Hu fitting formulae. Growth equations
were solved with adaptive Runge--Kutta methods
({\tt solve\_ivp} with {\tt RK45} or {\tt DOP853}).

Analytical derivations followed the DFD framework
established in Ref.~\cite{Alcock2025}, using the
spectral action formalism of
Refs.~\cite{Chamseddine2007,ConnesChamseddine2008}
for the $\alpha$-tower computations. All topological
results (Betti numbers, indices, Chern numbers) were
verified against standard references.

The campaign was conducted with Claude (Anthropic)
as a technical collaborator, providing calculation
verification, adversarial review, and independent
cross-checks. Each round's design was informed by
the previous round's results, with explicit attention
to closing identified gaps and testing proposed
mechanisms against all known constraints.

\subsection{Raw campaign outputs}

The individual analysis reports, Python solvers, and
synthesis documents from all 17~rounds are archived at
\url{https://doi.org/10.5281/zenodo.XXXXX} (to be
deposited upon submission). The archive contains
135+~files totaling ${\sim}4.3$~MB for Phase~I and
11~files totaling ${\sim}264$~KB for Phase~II, organized
by round. Each file is named with the convention
{\tt R\{round\}\_\{agent\}\_\{topic\}.\{ext\}}.

%==========================================================================
\section{Phase III: Verification and Tension Resolution
(37 follow-up analyses)}
\label{sec:phase3}
%==========================================================================

A third campaign of 37~analyses (25~broad + 9~targeted
+ 3~surgical) verified the companion paper against
2024--2026 observational data, resolved internal
tensions, and applied corrections.

\emph{(a)~Cosmic birefringence escalation.}
Combined Planck~PR4 + ACT~DR6 data yield
$\beta = 0.21^\circ\pm 0.03^\circ$. DFD predicts
$\beta=0^\circ$. All three DFD mechanisms (gravitational
CS from KK reduction, $\psi$-screen gradient,
$\CP^2$ geometric parity) confirmed unable to produce
photon birefringence. Galactic dust systematics remain
viable: measured $\beta$ decreases with larger masks,
ACT shows frequency discrepancy.

\emph{(b)~Sphaleron tension resolved.}
Two different $\Sph^3$'s: the internal
$\Sph^3\subset M_7$ protects $B{-}L$ via
$\pi_3(\Sph^3)=\mathbb{Z}$; the gauge
$\mathrm{SU}(2)_L$ configuration space has sphaleron
saddle points violating $B{+}L$ while preserving
$B{-}L$. The $16/3$ counting is safe.

\emph{(c)~$\psi$-screen reinterpretation.}
The $\chi$~field handles all acoustic physics;
the $\psi$-screen handles only the low-$z$ distance
bias ($\Delta\psi\approx 0.07$) that produces the
Hubble tension. The $c$-variation cancellation on the
sound horizon is a consistency feature.

\emph{(d)~Parameter count corrected} to three
($\Omega_b$, $n_s$, $A_s$), not zero.

\emph{(e)~Confirmed predictions:}
$H_0 = 72.09$ matches CCHP ($72.04$) to $0.07\%$;
$\sigma_8 = 0.811$ matches KiDS-Legacy ($0.815$) at
$0.2\sigma$; moduli stabilization $R/r = 1/\sqrt{3}$
derived from spectral action; Starobinsky $R^2$
inflation natural from the spectral action.

%==========================================================================
\section*{Acknowledgments}
%==========================================================================

The author thanks Claude (Anthropic) for technical
collaboration throughout the campaign.

\begin{thebibliography}{99}

\bibitem{Alcock2025}
G.~Alcock, ``Density Field Dynamics: A Scalar Refractive
Theory of Gravity---Unified Review v3.3,'' Zenodo (2025),
\href{https://doi.org/10.5281/zenodo.19029160}%
{10.5281/zenodo.19029160}.

\bibitem{Chamseddine2007}
A.~H.~Chamseddine, A.~Connes, and M.~Marcolli,
``Gravity and the Standard Model with neutrino mixing,''
Adv.\ Theor.\ Math.\ Phys.\ \textbf{11}, 991 (2007).

\bibitem{ConnesChamseddine2008}
A.~Connes and A.~H.~Chamseddine, ``Why the Standard Model,''
J.\ Geom.\ Phys.\ \textbf{58}, 38 (2008).

\end{thebibliography}

\end{document}
